\chapter{Exact Matching, Dictionaries, and String Indexes}
\label{ch:string-matching-indexes}

A scanner searching for a long delimiter can repeatedly recompare the same prefix after every mismatch.
An identifier dictionary can repeatedly compare the same initial letters across thousands of words.
Both problems disappear only when the implementation records the structure of the comparisons it has already performed.
Prefix functions, tries, and suffix indexes record different structure for different query workloads.

\section{Recover from mismatch without restarting}
\label{sec:ch78-prefix}
Let pattern $p$ have length $m$ and text $s$ length $n$, both measured in the same alphabet.
After matching $p[0:j]$, a mismatch does not erase the fact that a suffix of the matched region might be another pattern prefix.
A \term{border} is a proper prefix that is also a suffix.
The prefix function $\pi[i]$ is the length of the longest border of $p[0:i+1]$.

For \code{ababa}, the prefix values are $[0,0,1,2,3]$.
At the last position, \code{aba} is both prefix and suffix.
If a full match ends there, retaining border length three permits an overlapping match two positions later.
Table~\ref{tab:ch78-prefix} works every position, and Figure~\ref{fig:ch78-border} marks the reused region.

\begin{table}[htbp]
\caption{The prefix function records the largest already justified fallback.}
\label{tab:ch78-prefix}
\centering
\begin{tabular}{rrlr}
\toprule
Index & Symbol & Longest border & Length \\
\midrule
0 & a & empty & 0 \\
1 & b & empty & 0 \\
2 & a & a & 1 \\
3 & b & ab & 2 \\
4 & a & aba & 3 \\
\bottomrule
\end{tabular}
\end{table}

\begin{figure}[htbp]
\centering
\begin{tikzpicture}
\node[basic block] (p) at (0,0) {a b a b a};
\node[basic block] (b) at (0,-1.4) {reusable border: a b a};
\draw[cfg edge] (p)--node[right,font=\small]{fallback length 3}(b);
\end{tikzpicture}
\caption{A complete match retains a nonempty border when overlapping matches are possible.}
\label{fig:ch78-border}
\end{figure}

Algorithm~\ref{alg:ch78-prefix} constructs the failure information.
Each successful extension increments $j$ by one and each fallback decreases it.
Across the whole scan there are only $O(m)$ increases and decreases, giving $O(m)$ time and $O(m)$ storage under constant-time symbol comparison.
The while loop can be long at one position without making total preprocessing quadratic.

\begin{algorithm}[htbp]
\caption{Construct the prefix function of a nonempty pattern.}
\label{alg:ch78-prefix}
\begin{pseudocode}
\Require pattern $p$ of length $m>0$
\Ensure longest proper-border lengths $\pi[0:m]$
\State initialize $\pi$ to zero; $j\gets0$
\For{$i=1$ to $m-1$}
\While{$j>0$ and $p[i]\ne p[j]$}
\State $j\gets\pi[j-1]$
\EndWhile
\If{$p[i]=p[j]$} \State $j\gets j+1$ \EndIf
\State $\pi[i]\gets j$
\EndFor
\end{pseudocode}
\end{algorithm}

\section{Search with a maintained prefix invariant}
\label{sec:ch78-search}
Before processing the next text symbol, $j$ records the longest pattern prefix matching a suffix of the text consumed so far.
Fallback replaces that suffix with its longest feasible border; no text position moves backward.
Algorithm~\ref{alg:ch78-search} reports all matches, not merely the first.
After reporting a match it falls back instead of resetting to zero, which preserves overlap.

\begin{algorithm}[htbp]
\caption{Find every exact match, including overlapping matches.}
\label{alg:ch78-search}
\begin{pseudocode}
\Require text $s$, pattern $p$, its prefix function when nonempty
\Ensure all match start positions in increasing order
\If{$|p|=0$} \State return every boundary $0$ through $|s|$ \EndIf
\State $j\gets0$
\For{$i=0$ to $|s|-1$}
\While{$j>0$ and $s[i]\ne p[j]$} \State $j\gets\pi[j-1]$ \EndWhile
\If{$s[i]=p[j]$} \State $j\gets j+1$ \EndIf
\If{$j=|p|$}
\State report $i+1-|p|$; $j\gets\pi[j-1]$
\EndIf
\EndFor
\end{pseudocode}
\end{algorithm}

The total cost is $O(n+m+z)$ including preprocessing and $z$ reported positions.
Searching \code{ababababa} for \code{ababa} reports $0,2,4$.
The empty pattern policy reports all $n+1$ boundaries rather than dereferencing a nonexistent first pattern symbol.
The companion checks all 255 binary texts of length at most seven against all 31 binary patterns of length at most four, comparing with a direct start-position oracle.

A streaming implementation keeps only $j$, the prefix table, and a running text offset between chunks.
Resetting $j$ at chunk boundaries misses matches split across buffers.
If the alphabet is Unicode scalars, the decoder also retains incomplete byte sequences between chunks under its own error policy.
Matching state and decoding state must not be conflated.

\section{Index a dictionary with shared prefixes}
\label{sec:ch78-trie}
A \term{trie} assigns an edge to each next symbol and a final marker to each complete stored word.
Words \code{cat}, \code{cot}, and \code{cut} share their first edge but diverge at the vowel.
For a dictionary with $L$ total symbol occurrences, a plain trie has at most $L+1$ states.
Lookup follows at most $m$ edges for an $m$-symbol query, with transition costs determined by the child-map representation.

Figure~\ref{fig:ch78-trie} shows the branching needed by the spellchecker.
A final marker belongs to the state after the last symbol, not to an extra terminator character that users must supply.
A word can be both final and have outgoing edges, as when \code{compile} and \code{compiler} coexist.
Frequency and display metadata belong to final entries and need explicit merge rules under normalization.

\begin{figure}[htbp]
\centering
\begin{tikzpicture}[node distance=12mm and 15mm]
\node[basic block] (r) {root};\node[basic block,below=of r] (c) {c};
\node[basic block,below left=of c] (a) {ca};\node[basic block,below=of c] (o) {co};
\node[basic block,below right=of c] (u) {cu};
\node[basic block,below=of a] (at) {cat: final};
\node[basic block,below=of o] (ot) {cot: final};
\node[basic block,below=of u] (ut) {cut: final};
\draw[cfg edge] (r)--(c);\draw[cfg edge] (c)--(a);\draw[cfg edge] (c)--(o);\draw[cfg edge] (c)--(u);
\draw[cfg edge] (a)--(at);\draw[cfg edge] (o)--(ot);\draw[cfg edge] (u)--(ut);
\end{tikzpicture}
\caption{A trie shares prefixes while keeping each state's complete root path unique.}
\label{fig:ch78-trie}
\end{figure}

Suffix-equivalent trie states can be merged into a smaller acyclic automaton.
That preserves membership, but it destroys the unique-prefix property: one state can now be reached by several output strings.
This matters in approximate search.
Keeping one predecessor or one prefix per merged state can accidentally discard distinct candidate words even if membership tests still pass.
The case study deliberately keeps a trie and charges its extra storage rather than silently relying on an invalid enumeration invariant.

\section{Choose which side to preprocess}
\label{sec:ch78-indexes}
Prefix-function matching preprocesses the pattern; a trie preprocesses a dictionary.
A suffix array preprocesses a fixed text into the lexicographically ordered list of suffix starts.
Pattern matches form a contiguous interval in that order, so binary search can locate them.
With naive symbol-by-symbol comparisons, searching costs $O(m\log n+z)$ for pattern length $m$, text length $n$, and $z$ returned positions; more sophisticated longest-common-prefix structures reduce repeated comparisons.
Those bounds exclude suffix-array construction, whose cost depends on the construction algorithm and alphabet representation.

Table~\ref{tab:ch78-indexes} gives a workload map rather than declaring one universal winner.
A compiler compiling a single source snapshot often values streaming scans and small auxiliary state.
A source browser answering repeated substring queries may justify a larger persistent index.
Dynamic edits require invalidation or a genuinely dynamic indexing algorithm; rebuilding a static suffix array is not constant-time update.

\begin{table}[htbp]
\caption{The reusable side of the workload determines which index can amortize its cost.}
\label{tab:ch78-indexes}
\centering
\begin{tabular}{lll}
\toprule
Structure & Preprocessed data & Natural query \\
\midrule
prefix function & one pattern & streaming exact matches \\
trie & word dictionary & prefix lookup and completion \\
multi-pattern automaton & pattern collection & all patterns in one stream \\
suffix array & fixed text snapshot & repeated substring lookup \\
hash table & complete keys & exact dictionary equality \\
\bottomrule
\end{tabular}
\end{table}

Multi-pattern matching augments shared-prefix states with failure transitions and output lists so one scan can recognize many patterns.
The output volume can itself be large when patterns are nested suffixes, and must appear in the complexity bound.
Fingerprint-based search can accelerate candidate filtering, but a collision-aware implementation verifies candidates unless its specification explicitly permits a probabilistic answer.
For source transformations and diagnostics, a false positive is normally unacceptable.

\section{Exercises}
\label{sec:ch78-exercises}
\begin{enumerate}
\item\label{ex:ch78-prefix} Compute the prefix function of \code{aaaa} and all its matches in \code{aaaaaa}.
\item\label{ex:ch78-stream} Split \code{ababababa} after its fourth symbol.
Which state must cross the chunk boundary to retain all \code{ababa} matches?
\item\label{ex:ch78-trie} Insert \code{compile} and \code{compiler} into an initially empty trie.
Count states including the root, and identify the two final states.
\item\label{ex:ch78-merge} Explain how merging suffix-equivalent dictionary states can preserve membership while invalidating an enumeration algorithm storing one prefix per state.
\end{enumerate}

\section{Further reading}
\label{sec:ch78-reading}
Crochemore, Hancart, and Lecroq develop exact matching, automata representations, and indexing in detail~\cite{crochemore-2007-strings}.
The lexer-construction chapters provide complementary automaton-building techniques; Chapter~\ref{ch:fst-spellchecker} uses the trie as a constrained output language rather than an exact-match table.
